\paragraph{A common multiplicative disk at 2.}
Let $v(2)=1$. For any two cyclic order-$\ell$ subgroups $H,J$ of
$E_0$, with $\ell$ odd, we claim
\[
 v\big(j(E_0/H)/j(E_0/J)-1\big)\ge4.                 \tag{G0}
\]
All valuations may be taken in one finite extension of $\mathbb Q_2$.
The integral invariant $1728$ gives $E_0$ potentially good reduction.
Take a field of good reduction containing $E_0[3\ell]$. Write
$x=r+u^2X$ for the change from the original short equation to a good
integral Weierstrass equation with coordinate $X$. Its discriminant
is a unit, while the original discriminant is $64$, so $v(u^2)=1$.
Every nonzero odd-torsion point under consideration has integral $X$:
reduction is injective on prime-to-$2$ torsion, so its reduction is
an affine point of the good model. In particular for a fixed
nonzero $3$-torsion point $T$ and every nonzero $Q\in H\cup J$ we
have $v(x_Q-x_T)\ge1$.

The $3$-division polynomial gives $3x_T^4-6x_T^2-1=0$.
Its roots are units. For $y=5x_T^2-3$ this yields
$3y^2-12y-88=0$. The Newton polygon has its only slope $-3/2$,
so $v(y)=3/2$. Since $x_Q=x_T+2z$ with $z$ integral,
$v(x_Q^2-x_T^2)\ge2$. For $g(x)=x(5x^2-3)$ it follows that
\[
 v(g(x_T))=3/2,\qquad v(g(x_Q)-g(x_T))\ge2.
\]
V\'elu's quotient model is $Y^2=X^3+A_H X+B_H$, where
\[
 A_H=-1-10\!\sum_{Q\in(H\setminus\{0\})/\{\pm1\}}(3x_Q^2-1),
 \qquad
 B_H=-14\!\sum_{Q\in(H\setminus\{0\})/\{\pm1\}}g(x_Q).
\]
These are the short-equation specialization of
\href{https://ddd.uab.cat/pub/pubmat/pubmat_a2007vEXTRA/pubmat_a2007vextrap147.pdf}{Miret--Moreno--Rio,
\emph{Generalization of V\'elu's formulae for isogenies between elliptic curves},
pp.~149--151, especially Theorem~1}.
In particular $A_H$ is a unit congruent to $-1$ modulo $2$ and
$v(B_H)\ge5/2$. The $H$ and $J$ sums have the same number
$(\ell-1)/2$ of terms, all congruent to $g(x_T)$ modulo $4$ in
the second sum. Thus $v(B_H-B_J)\ge3$ and
$v(B_H+B_J)=v(B_H-B_J+2B_J)\ge3$.
This gives $v(B_H^2-B_J^2)\ge6$.

Put $s_H=27B_H^2/(4A_H^3)$. Then $v(s_H)\ge3$ and
$v(s_H-s_J)\ge4$: replacing $B_H^2$ by $B_J^2$ costs at least four
in valuation after division by $4$, while
$v(A_H^{-3}-A_J^{-3})\ge1$. The identity
$j(E_0/H)=1728/(1+s_H)$ proves (G0).
Fixing any one of the quotient invariants $j_*$ therefore puts all
$h$ conjugates in $j_*(1+16\mathcal O)$, where $\mathcal O$ is the
valuation ring. The center is allowed to be algebraic. The recurrence
and the root count below apply over this complete finite extension
with the same normalized valuation.

The depth four is sharp when the disk assertion ranges over all odd
prime conductors. For $\ell=3$ let $Y=j/1728-1$. Substitution of
$A=9-30x_T^2$, $B=-14x_T(5x_T^2-3)$ into the formula for $j$
and elimination by $3x_T^4-6x_T^2-1=0$ give the quadratic
$27Y^2-2399040Y-18593344$. Its discriminant has $2$-adic valuation
eight. The two roots differ by an element of valuation four, and
both $1+Y$ are units. The corresponding pair thus attains equality
in (G0). The archive reproduces the elimination and also checks
the corresponding degree-five isogenies in characteristic zero.
